\section{October 2, 2026}

\subsection{Operators and invariant subspaces}

\begin{definition}[Linear operator]
  A linear map from a vector space $V$ to itself is called a
  \emph{linear operator} on $V$. We write
  $\mathcal L(V)=\mathcal L(V,V)$ for the space of all operators
  on $V$.
\end{definition}

One goal in studying an operator $T\in\mathcal L(V)$ is to find
a direct sum decomposition
\[
  V=V_1\oplus\cdots\oplus V_n
\]
into subspaces that $T$ preserves. We can then study $T$ by
studying its restrictions to these subspaces.

\begin{definition}[Invariant subspace]
  Let $T\in\mathcal L(V)$. A subspace $U\subseteq V$ is
  \emph{invariant under $T$} if $T(u)\in U$ for every $u\in U$;
  equivalently, $T(U)\subseteq U$.
\end{definition}

When $U$ is invariant under $T$, the restriction
$T|_U\colon U\to U$ is itself an operator on $U$.

\begin{example}
  Define $D\in\mathcal L(\mathcal P(\bb{R}))$ by $D(p)=p''$.
  The subspace $\mathcal P_5(\bb{R})$ is invariant under $D$:
  taking two derivatives of a polynomial of degree at most $5$
  produces a polynomial of degree at most $3$, or the zero
  polynomial, and hence an element of $\mathcal P_5(\bb{R})$.
\end{example}

\begin{example}
  For every $T\in\mathcal L(V)$, the subspaces
  \[
    \{0\},\qquad V,\qquad \ker T,\qquad \operatorname{im}T
  \]
  are invariant under $T$. The first two assertions follow from
  $T(0)=0$ and the fact that $T$ maps $V$ to itself. If
  $u\in\ker T$, then $T(u)=0\in\ker T$. If
  $u\in\operatorname{im}T$, then $u=T(v)$ for some $v\in V$,
  and $T(u)=T(T(v))\in\operatorname{im}T$.
\end{example}

\subsection{Eigenvalues and eigenvectors}

Consider a one-dimensional subspace of $V$. Let
$v\in V\setminus\{0\}$ and $U=\operatorname{span}\{v\}$. If
$U$ is invariant under $T$, then $T(v)\in U$, so there is a
scalar $\lambda\in\bb{F}$ such that
\[
  T(v)=\lambda v.
\]
Conversely, if $T(v)=\lambda v$, then for every $a\in\bb{F}$,
\[
  T(av)=aT(v)=a\lambda v\in U.
\]
Thus $U$ is invariant under $T$ precisely when $T(v)$ is a
scalar multiple of $v$.

\begin{definition}[Eigenvalue and eigenvector]
  Let $T\in\mathcal L(V)$. A scalar $\lambda\in\bb{F}$ is an
  \emph{eigenvalue} of $T$ if there exists a nonzero vector
  $v\in V$ such that $T(v)=\lambda v$. Such a vector $v$ is
  called an \emph{eigenvector} of $T$ corresponding to $\lambda$.
\end{definition}

Consequently, $V$ has a one-dimensional subspace invariant under
$T$ if and only if $T$ has an eigenvalue.

\begin{proposition}\label{prop:eigenvalue-equivalences}
  Suppose that $V$ is finite-dimensional, $T\in\mathcal L(V)$,
  and $\lambda\in\bb{F}$. The following are equivalent:
  \begin{enumerate}[label=(\roman*)]
    \item $\lambda$ is an eigenvalue of $T$;
    \item $T-\lambda\operatorname{Id}_V$ is not injective;
    \item $T-\lambda\operatorname{Id}_V$ is not surjective;
    \item $T-\lambda\operatorname{Id}_V$ is not invertible.
  \end{enumerate}
\end{proposition}

\begin{proof}
  For $v\in V$,
  \[
    T(v)=\lambda v
      \quad\Longleftrightarrow\quad
      (T-\lambda\operatorname{Id}_V)(v)=0.
  \]
  Hence $\lambda$ is an eigenvalue if and only if
  $\ker(T-\lambda\operatorname{Id}_V)\neq\{0\}$, which is
  equivalent to $T-\lambda\operatorname{Id}_V$ not being
  injective. This proves (i)$\Longleftrightarrow$(ii).
  Since $T-\lambda\operatorname{Id}_V$ is an operator on a
  finite-dimensional space,
  Proposition~\ref{prop:equal-dimension-invertibility} gives the
  equivalence of (ii), (iii), and (iv).
\end{proof}

In particular, the eigenvectors corresponding to $\lambda$ are
exactly the nonzero vectors in
$\ker(T-\lambda\operatorname{Id}_V)$. This kernel includes the
zero vector, but the zero vector is never an eigenvector.

\subsection{A rotation over the real and complex numbers}

\begin{example}\label{ex:rotation-eigenvalues}
  Let $T\in\mathcal L(\bb{F}^2)$ be defined by
  \[
    T\begin{bmatrix}w\\z\end{bmatrix}
      =\begin{bmatrix}-z\\w\end{bmatrix}.
  \]
  Its matrix with respect to the standard basis is
  \[
    A=\begin{bmatrix}0&-1\\1&0\end{bmatrix}.
  \]

  If $\bb{F}=\bb{R}$, then $T$ is counterclockwise rotation by
  $90^\circ$. No nonzero real vector is taken to a scalar
  multiple of itself. Thus $T$ has no real eigenvalues and no
  eigenvectors in $\bb{R}^2$.

  To see algebraically how the field matters, the eigenvector
  equation is
  \[
    \begin{bmatrix}-z\\w\end{bmatrix}
      =\lambda\begin{bmatrix}w\\z\end{bmatrix},
    \qquad\text{so}\qquad
    \lambda w=-z,\quad \lambda z=w.
  \]
  If $z=0$, then $w=0$, which is excluded for an eigenvector.
  Therefore $z\neq0$, and substitution gives
  \[
    \lambda^2z=-z
      \quad\Longrightarrow\quad \lambda^2=-1.
  \]
  There is no solution in $\bb{R}$. Over $\bb{C}$, the only
  possibilities are $\lambda=i$ and $\lambda=-i$, and both
  occur:
  \begin{align*}
    T\begin{bmatrix}1\\-i\end{bmatrix}
      &=\begin{bmatrix}i\\1\end{bmatrix}
        =i\begin{bmatrix}1\\-i\end{bmatrix},\\
    T\begin{bmatrix}1\\i\end{bmatrix}
      &=\begin{bmatrix}-i\\1\end{bmatrix}
        =-i\begin{bmatrix}1\\i\end{bmatrix}.
  \end{align*}
  More precisely,
  \[
    \ker(T-i\operatorname{Id}_{\bb{C}^2})
      =\operatorname{span}\left\{\begin{bmatrix}1\\-i\end{bmatrix}\right\},
    \qquad
    \ker(T+i\operatorname{Id}_{\bb{C}^2})
      =\operatorname{span}\left\{\begin{bmatrix}1\\i\end{bmatrix}\right\}.
  \]
  Thus the eigenvectors for $i$ are $(w,-iw)$ with $w\neq0$,
  and those for $-i$ are $(w,iw)$ with $w\neq0$.
\end{example}

\subsection{Eigenvectors corresponding to distinct eigenvalues}

\begin{proposition}\label{prop:distinct-eigenvalues-independent}
  Let $T\in\mathcal L(V)$. If $v_1,\ldots,v_n$ are eigenvectors
  of $T$ corresponding to distinct eigenvalues
  $\lambda_1,\ldots,\lambda_n$, then $v_1,\ldots,v_n$ are
  linearly independent.
\end{proposition}

\begin{proof}
  Suppose, to the contrary, that the list is linearly dependent.
  Choose a dependent sublist of smallest length $m$ and relabel
  it as $v_1,\ldots,v_m$. Denote their eigenvalues by
  $\lambda_1,\ldots,\lambda_m$. All eigenvectors are nonzero,
  so $m\geq2$. There are scalars $a_1,\ldots,a_m$, not
  all zero, such that
  \begin{equation}\label{eq:eigenvector-dependence}
    a_1v_1+\cdots+a_mv_m=0.
  \end{equation}
  Minimality of $m$ implies that every $a_j$ is nonzero;
  otherwise, omitting a term would give a shorter dependent
  sublist. Applying $T$ to
  \eqref{eq:eigenvector-dependence} gives
  \[
    a_1\lambda_1v_1+\cdots+a_m\lambda_mv_m=0.
  \]
  Subtracting $\lambda_1$ times
  \eqref{eq:eigenvector-dependence}, we obtain
  \[
    a_2(\lambda_2-\lambda_1)v_2
      +\cdots+a_m(\lambda_m-\lambda_1)v_m=0.
  \]
  Each coefficient is nonzero because $a_j\neq0$ and the
  eigenvalues are distinct. This is a nontrivial dependence
  relation among $m-1$ of the eigenvectors, contradicting the
  minimality of $m$.
\end{proof}

\begin{corollary}\label{cor:eigenvalue-count-bound}
  Every operator on a finite-dimensional vector space $V$ has
  at most $\dim V$ distinct eigenvalues.
\end{corollary}

\begin{proof}
  Choose one eigenvector for each distinct eigenvalue. By
  Proposition~\ref{prop:distinct-eigenvalues-independent}, any
  finite list of these vectors is linearly independent. Such a
  list has length at most $\dim V$, so there cannot be more
  than $\dim V$ distinct eigenvalues.
\end{proof}
